% Auto-generated by tables/scripts/make_rep_quality.py --- do not hand-edit.
% Single late checkpoint at 1000k (latest step all four share at 3 seeds).
\begin{table}[tbp]
\centering
\small
\setlength{\tabcolsep}{5pt}
\begin{tabular}{lccccc}
\toprule
 & & & \multicolumn{3}{c}{\textbf{CATH top-1 acc.\,$\uparrow$}} \\
\cmidrule(lr){4-6}
\textbf{Variant} & \textbf{IF top-1\,$\uparrow$} & \textbf{dihedral MAE\,(\,$^\circ$\,)\,$\downarrow$} & C & A & T \\
\midrule
Baseline & $0.129_{\pm0.004}$ & $31.5_{\pm1.7}$ & $0.724_{\pm0.013}$ & $0.408_{\pm0.007}$ & $0.332_{\pm0.009}$ \\
REPA L4-random & \rd{$-0.010_{\pm0.002}$} & $+3.9_{\pm3.1}$ & \gd{$+0.049_{\pm0.005}$} & $+0.011_{\pm0.008}$ & $-0.005_{\pm0.013}$ \\
REPA L9-GearNet & \gd{$+0.023_{\pm0.002}$} & \gb{$-6.4_{\pm0.3}$} & \gb{$+0.153_{\pm0.009}$} & \gb{$+0.267_{\pm0.006}$} & \gb{$+0.415_{\pm0.017}$} \\
REPA L9-MPNN & \gb{$+0.026_{\pm0.002}$} & $-5.6_{\pm4.0}$ & \gd{$+0.081_{\pm0.006}$} & \gd{$+0.063_{\pm0.003}$} & \gd{$+0.106_{\pm0.026}$} \\
\bottomrule
\end{tabular}
\caption{\textbf{REPA offers persistent and encoder-routed representation advantages.} At a representative late checkpoint, REPA-GearNet wins the fold probes (CATH C/A/T) decisively, while REPA-MPNN edges the per-residue probe (IF). The random control gains least and does not help the per-residue probes, evidence the effect is genuine structural transfer rather than regularisation. We headline on CATH-A and report C and T for completeness. (Single checkpoint at 1.0M --- the latest training step all four variants share at three seeds, and a representative point on the baseline trajectory rather than a transient swing; $n{=}3$ seeds, with $\pm$ half the seed range. Best-layer linear probe at $n_\text{train}{=}1000$ (cross-database set for IF/dihedral, cleantrain for CATH); dihedral carries the widest seed spread, so the per-residue read rests on IF. Colour marks significance against the combined seed spread: green better (darker $=$ column-best), red worse, uncoloured within seed noise. Persistence across training is shown in Figure~\ref{fig:proteina-rep}.)}
\label{tab:proteina-rep}
\end{table}
